\section{Programmieraufgaben und generierte Testfälle}
\label{anhang:programmieraufgaben-testfaelle}

\lstdefinestyle{appendixdata}{
    basicstyle=\ttfamily\footnotesize,
    breaklines=true,
    breakatwhitespace=false,
    columns=fullflexible,
    keepspaces=true,
    showstringspaces=false,
    frame=single,
    framerule=0.4pt,
    numbers=none,
    tabsize=2
}
Die folgende Auflisting enthält die Aufgabenbeschreibungen und Eingabebeschränkungen (Constraints), die daraus generierten Testfälle 
und eine Erklärung der jeweiligen Aufgabe. 

\subsection{Problem 1: Binary Number with Alternating Bits (Leetcode 693)}

\subsubsection*{Aufgabenbeschreibung}
Given a positive integer, check whether it has alternating bits: namely, if two adjacent bits will always have different values

\subsubsection*{Constraints}
\begin{lstlisting}[style=appendixdata]
1 <= n <= 2^31 - 1
\end{lstlisting} 

\subsubsection*{Generierte Testfälle}
\begin{lstlisting}[style=appendixdata]
{
  "test1": {"n": 5, "expected": true},
  "test2": {"n": 7, "expected": false},
  "test3": {"n": 10, "expected": true},
  "test4": {"n": 11, "expected": false},
  "test5": {"n": 1, "expected": true},
  "test6": {"n": 0, "expected": false},
  "test7": {"n": 1431655765, "expected": true},
  "test8": {"n": 2863311530, "expected": false},
  "test9": {"n": 2, "expected": true},
  "test10": {"n": 3, "expected": false},
  "test11": {"n": 170, "expected": true},
  "test12": {"n": 85, "expected": true},
  "test13": {"n": 2147483647, "expected": false},
  "test14": {"n": 1431655766, "expected": false},
  "test15": {"n": 2147483646, "expected": false}
}
\end{lstlisting}

\subsubsection*{Erklärung der Aufgabe}
Die Binärdarstellung einer Positiven Ganzzahl soll darauf geprüft werden, 
ob sich benachbarte Bits überall unterscheiden. Algorithmisch muss die Bitfolge von rechts nach links durchgegangen werden um jedes Bit 
mit seinem direkten Nachbarn zu vergleichen.

\clearpage

\subsection{Problem 2: Sort Vowels by frequrency (Leetcode 3913)}

\subsubsection*{Aufgabenbeschreibung}
You are given a string s consisting of lowercase English characters.
Rearrange only the vowels in the string so that they appear in non-increasing order of their frequency.
If multiple vowels have the same frequency, order them by the position of their first occurrence in s.
Return the modified string.
Vowels are 'a', 'e', 'i', 'o', and 'u'.
The frequency of a letter is the number of times it occurs in the string.

Example:

Input: s = 'leetcode'

Output: 'leetcedo'

Explanation: Vowels in the string are ['e', 'e', 'o', 'e'] with frequencies: e = 3, o = 1. Sorting in non-increasing order of frequency and placing them back 
into the vowel positions results in "leetcedo".

\subsubsection*{Constraints}
\begin{lstlisting}[style=appendixdata]
1 <= s.length <= 10^5
s consists of lowercase English letters
\end{lstlisting}

\subsubsection*{Generierte Testfälle}
\begin{lstlisting}[style=appendixdata]
{
  "test1": {"s": "leetcode", "expected": "leetcedo"},
  "test2": {"s": "aeiou", "expected": "aeiou"},
  "test3": {"s": "aabbccddeeiioouu", "expected": "aabbccddeeiioouu"},
  "test4": {"s": "uauaueiouaeiou", "expected": "uuaaaeeiiouou"},
  "test5": {"s": "xyz", "expected": "xyz"},
  "test6": {"s": "aebecidofu", "expected": "aaebecidofu"},
  "test7": {"s": "oooeeeiiiaaa", "expected": "oooeeeiiiaaa"},
  "test8": {"s": "aeaeaeaeaeaeae", "expected": "aeaeaeaeaeaeae"},
  "test9": {"s": "aeiouaeiouaeiou", "expected": "aeiouaeiouaeiou"},
  "test10": {"s": "bbeeiioouu", "expected": "bbeeiiouou"},
  "test11": {"s": "aeaeiouioiuoaeae", "expected": "aeaeaeaeouiouiou"},
  "test12": {"s": "azbzczdz", "expected": "azbzczdz"},
  "test13": {"s": "eioaeuu", "expected": "uuueioa"},
  "test14": {"s": "aaeeiioouu", "expected": "aaeeiioouu"},
  "test15": {"s": "uuuuiiiiiaaaaeeee", "expected": "uuuuiiiiiaaaaeeee"}
}
\end{lstlisting}

\subsubsection*{Erklärung der Aufgabe}
Innerhalb einer Zeichenkette sollen ausschließlich die Vokale nach absteigender Häufigkeit umsortiert werden. Alle Konsonanten 
sollen ihre ursprüngliche Postition beibehalten. Algorithmisch besteht die Herausforderung darin, dass die Vokalhäufigkeiten und die intiale Ordnung der Vokale bestimmt 
werden muss, woraufhin eine sortierte Vokalfolge erstellt werden kann um diese auf die ursprüngliche Zeichenkette anzuwenden.

\clearpage

\subsection{Problem 3: Rotate List (Leetcode 61)}

\subsubsection*{Aufgabenbeschreibung}
Given the head of a linked list, rotate the list to the right by k places.

\subsubsection*{Constraints}
\begin{lstlisting}[style=appendixdata]
The number of nodes in the list is in the range [0, 500].
-100 <= Node.val <= 100
0 <= k <= 2 * 10^9
\end{lstlisting}

\subsubsection*{Generierte Testfälle}
\begin{lstlisting}[style=appendixdata]
{
  "test1": {"head": [], "k": 0, "expected": []},
  "test2": {"head": [], "k": 5, "expected": []},
  "test3": {"head": [1], "k": 0, "expected": [1]},
  "test4": {"head": [1], "k": 1, "expected": [1]},
  "test5": {"head": [1], "k": 1000000000, "expected": [1]},
  "test6": {"head": [1,2,3,4,5], "k": 0, "expected": [1,2,3,4,5]},
  "test7": {"head": [1,2,3,4,5], "k": 1, "expected": [5,1,2,3,4]},
  "test8": {"head": [1,2,3,4,5], "k": 2, "expected": [4,5,1,2,3]},
  "test9": {"head": [1,2,3,4,5], "k": 5, "expected": [1,2,3,4,5]},
  "test10": {"head": [1,2,3,4,5], "k": 7, "expected": [4,5,1,2,3]},
  "test11": {"head": [1,2,3,4,5], "k": 13, "expected": [3,4,5,1,2]},
  "test12": {"head": [10,-10,20,-20,30,-30], "k": 3, "expected": [20,-20,30,-30,10,-10]},
  "test13": {"head": [0,0,0,0], "k": 2, "expected": [0,0,0,0]},
  "test14": {"head": [100,99,-99,-100], "k": 4, "expected": [100,99,-99,-100]},
  "test15": {"head": [1,2], "k": 1000000001, "expected": [2,1]}
}
\end{lstlisting}

\subsubsection*{Erklärung der Aufgabe}
Eine verkette Liste soll um eine vorgegebene Anzahl an Postitionen (k) nach rechts verschoben werden. Das Verschieben der einzelnen 
Knoten wird durch das Verbinden der verketten Liste und das Trennen an einer anderen Stelle abgebildet.

\clearpage

\subsection{Problem 4: Find Highest Altitude(Leetcode 1732)}

\subsubsection*{Aufgabenbeschreibung}
There is a biker going on a road trip. The road trip consists of n + 1 points at various altitudes. The biker starts his trip on point 0 with altitude equal 0.
You are given an integer array gain of length n where gain[i] is the net gain in altitude between points i and i + 1 for all (0 <= i < n). Return the highest altitude of a point.

\subsubsection*{Constraints}
\begin{lstlisting}[style=appendixdata]
n == gain.length
1 <= n <= 100
-100 <= gain[i] <= 100
\end{lstlisting}

\subsubsection*{Generierte Testfälle}
\begin{lstlisting}[style=appendixdata]
{
  "test1": { "gain": [-5, 1, 5, 0, -7], "expected": 1 },
  "test2": { "gain": [0, 0, 0, 0], "expected": 0 },
  "test3": { "gain": [100, 100, 100], "expected": 300 },
  "test4": { "gain": [-100, -100, -100], "expected": 0 },
  "test5": { "gain": [100, -50, -40, 10], "expected": 100 },
  "test6": { "gain": [-2, -1, -2, -1], "expected": 0 },
  "test7": { "gain": [1, 2, 3, 4, 5], "expected": 15 },
  "test8": { "gain": [-1, -2, -3, -4, -5], "expected": 0 },
  "test9": { "gain": [10, -5, 15, -10, 20], "expected": 30 },
  "test10": { "gain": [-10, 5, -15, 10, -20], "expected": 0 },
  "test11": { "gain": [0], "expected": 0 },
  "test12": { "gain": [100], "expected": 100 },
  "test13": { "gain": [-100], "expected": 0 },
  "test14": { "gain": [1, -1, 1, -1, 1, -1], "expected": 1 },
  "test15": { "gain": [50, -25, 25, -50, 50], "expected": 50 }
}
\end{lstlisting}

\subsubsection*{Erklärung der Aufgabe}
Die Aufgabe beschreibt eine Autofahrt welche auf Höhe 0 beginnt. Ein Array beinhaltet die relativen Höhenänderungen zwischen 
aufeinanderfolgenden Punkten. Gesucht wird die absolut höchste Stelle, was algorithmisch der Bestimmung der höchsten Präfixsumme entspricht.

\clearpage

\subsection{Problem 5: Minimum Flips (Leetcode 1318)}

\subsubsection*{Aufgabenbeschreibung}
Given 3 positives numbers a, b and c. Return the minimum flips required in some bits of a and b to make ( a OR b == c ). (bitwise OR operation). Flip operation consists of change any single bit 1 to 0 or change the bit 0 to 1 in their binary representation.

\subsubsection*{Constraints}
\begin{lstlisting}[style=appendixdata]
1 <= a <= 10^9
1 <= b <= 10^9
1 <= c <= 10^9
\end{lstlisting}

\subsubsection*{Generierte Testfälle}
\begin{lstlisting}[style=appendixdata]
{
  "test1": {"a": 0, "b": 0, "c": 0, "expected": 0},
  "test2": {"a": 0, "b": 0, "c": 1, "expected": 1},
  "test3": {"a": 1, "b": 1, "c": 0, "expected": 2},
  "test4": {"a": 2, "b": 4, "c": 6, "expected": 0},
  "test5": {"a": 2, "b": 4, "c": 7, "expected": 1},
  "test6": {"a": 8, "b": 7, "c": 15, "expected": 0},
  "test7": {"a": 8, "b": 7, "c": 0, "expected": 5},
  "test8": {"a": 15, "b": 15, "c": 0, "expected": 8},
  "test9": {"a": 10, "b": 5, "c": 15, "expected": 0},
  "test10": {"a": 10, "b": 5, "c": 0, "expected": 6},
  "test11": {"a": 7, "b": 0, "c": 15, "expected": 1},
  "test12": {"a": 1, "b": 2, "c": 4, "expected": 3},
  "test13": {"a": 3, "b": 5, "c": 7, "expected": 0},
  "test14": {"a": 1, "b": 1, "c": 3, "expected": 1},
  "test15": {"a": 31, "b": 0, "c": 16, "expected": 4}
}
\end{lstlisting}

\subsubsection*{Erklärung der Aufgabe}
Für drei positive Ganzzahlen (a, b, c) größer Null soll die minimale Anzahl an Bitänderungen in a und b ermittelt 
werden, sodass ihr bitweises logisches Oder dem Wert c entspricht. Das Problem kann gelöst werden, indem Bit für Bit geprüft wird, ob ein Bit gesetzt oder gelöscht werden muss 
und ob eine oder zwei Änderungen dafür erforderlich sind.

\clearpage

\subsection{Problem 6: Sort Colors (Leetcode 75)}

\subsubsection*{Aufgabenbeschreibung}
Given an array nums with n objects colored red, white, or blue, sort them in-place so that objects of the same color are adjacent, with the colors in the order red, white, and blue.
We will use the integers 0, 1, and 2 to represent the color red, white, and blue, respectively.
You must solve this problem without using the library's sort function.

\subsubsection*{Constraints}
\begin{lstlisting}[style=appendixdata]
n == nums.length
1 <= n <= 300
nums[i] is either 0, 1, or 2.
\end{lstlisting}

\subsubsection*{Generierte Testfälle}
\begin{lstlisting}[style=appendixdata]
{
  "test1": {"nums": [0], "expected": [0]},
  "test2": {"nums": [2], "expected": [2]},
  "test3": {"nums": [1], "expected": [1]},
  "test4": {"nums": [0, 1, 2], "expected": [0, 1, 2]},
  "test5": {"nums": [2, 1, 0], "expected": [0, 1, 2]},
  "test6": {"nums": [0, 0, 0, 0], "expected": [0, 0, 0, 0]},
  "test7": {"nums": [2, 2, 2, 2], "expected": [2, 2, 2, 2]},
  "test8": {"nums": [1, 1, 1, 1], "expected": [1, 1, 1, 1]},
  "test9": {"nums": [0, 2, 1, 0, 2, 1], "expected": [0, 0, 1, 1, 2, 2]},
  "test10": {"nums": [2, 2, 1, 0, 0, 1, 1, 2, 0], "expected": [0, 0, 0, 1, 1, 1, 2, 2, 2]},
  "test11": {"nums": [1, 0, 0, 2, 1, 2, 0], "expected": [0, 0, 0, 1, 1, 2, 2]},
  "test12": {"nums": [1, 2, 1, 0, 0, 2, 1, 0, 2], "expected": [0, 0, 0, 1, 1, 1, 2, 2, 2]},
  "test13": {"nums": [0, 1], "expected": [0, 1]},
  "test14": {"nums": [2, 0], "expected": [0, 2]},
  "test15": {"nums": [1, 2], "expected": [1, 2]}
}
\end{lstlisting}

\subsubsection*{Erklärung der Aufgabe}
Ein Array aus den Werten 0, 1 und 2, welche die Farben rot, weiß und blau representieren, soll ohne die Verwendung von Bibliotheksfunktionen 
in-place in aufsteigende Reihenfolge sortiert werden.

\clearpage

\subsection{Problem 7: Max Number of K-Sum Pairs (Leetcode 1679)}

\subsubsection*{Aufgabenbeschreibung}
You are given an integer array nums and an integer k.
In one operation, you can pick two numbers from the array whose sum equals k and remove them from the array.
Return the maximum number of operations you can perform on the array.

\subsubsection*{Constraints}
\begin{lstlisting}[style=appendixdata]
1 <= nums.length <= 10^5
1 <= nums[i] <= 10^9
1 <= k <= 10^9
\end{lstlisting}

\subsubsection*{Generierte Testfälle}
\begin{lstlisting}[style=appendixdata]
{
  "test1": {"nums": [1, 2, 3, 4], "k": 5, "expected": 2},
  "test2": {"nums": [3, 1, 3, 4, 3], "k": 6, "expected": 1},
  "test3": {"nums": [1], "k": 2, "expected": 0},
  "test4": {"nums": [2, 2, 2, 2], "k": 4, "expected": 2},
  "test5": {"nums": [5, 5, 5, 5, 5], "k": 10, "expected": 2},
  "test6": {"nums": [1, 9, 2, 8, 3, 7, 4, 6, 5], "k": 10, "expected": 4},
  "test7": {"nums": [1, 2, 3, 4, 5], "k": 10, "expected": 0},
  "test8": {"nums": [1000000000, 1], "k": 1000000001, "expected": 1},
  "test9": {"nums": [1, 1, 1, 1, 1, 1], "k": 2, "expected": 3},
  "test10": {"nums": [4, 4, 4, 4, 4], "k": 8, "expected": 2},
  "test11": {"nums": [1, 2, 3, 7, 8, 9], "k": 10, "expected": 3},
  "test12": {"nums": [5, 5, 5, 5, 5, 5, 5], "k": 10, "expected": 3},
  "test13": {"nums": [1, 3, 5, 7, 9], "k": 8, "expected": 2},
  "test14": {"nums": [1], "k": 1, "expected": 0},
  "test15": {"nums": [1, 1, 2, 2, 3, 3, 4, 4], "k": 5, "expected": 4}
}
\end{lstlisting}

\subsubsection*{Erklärung der Aufgabe}
Aus einem Zahlen-Array sollen möglichst viele Paare gefunden und entfernt werden, deren Summe einem vorgegebenen Wert k entspricht. 
Die algorithmische Herausforderung besteht darin, passende Komplementwerte unter Berücksichtigung ihrer Häufigkeiten zuzuordnen, sodass jedes Element höchstens einmal 
verwendet wird.
